% ===== PRÉAMBULE CANONIQUE — à coller VERBATIM en tête de CHAQUE .tex =====
\documentclass[11pt,a4paper]{article}
\usepackage[utf8]{inputenc}\usepackage[T1]{fontenc}\usepackage{lmodern}
\usepackage[french]{babel}
\usepackage{amsmath,amssymb,mathtools}\usepackage{icomma}
\usepackage[version=4]{mhchem}          % \ce{...} : équations & formules
\usepackage{siunitx}\sisetup{locale=FR} % \qty{}{}, \si{}, \ang{}
\usepackage{chemfig}                    % \chemfig{...} : structures développées
\usepackage{xcolor}\usepackage{colortbl}
\usepackage{tikz}\usepackage{pgfplots}\pgfplotsset{compat=1.18}
\usepackage[most]{tcolorbox}\usepackage{enumitem}\usepackage{array,booktabs}
\usepackage[a4paper,margin=2.2cm]{geometry}
\usetikzlibrary{arrows.meta,calc,angles,quotes,positioning,patterns,backgrounds,shapes.geometric}
\definecolor{primary}{RGB}{222,49,99}\definecolor{primarydark}{RGB}{150,22,55}
\definecolor{defcol}{RGB}{0,114,178}\definecolor{methcol}{RGB}{52,92,148}
\definecolor{excol}{RGB}{0,135,90}\definecolor{attncol}{RGB}{200,40,55}
\tcbset{boxbase/.style={enhanced,breakable,boxrule=0.9pt,arc=2.5pt,
  left=3mm,right=3mm,top=2.3mm,bottom=2.3mm,fonttitle=\bfseries,coltitle=white}}
% --- boîtes TUTORIEL (noms EXACTS, ne pas renommer) ---
\newtcolorbox{cle}[1][]{boxbase,colback=defcol!6,colframe=defcol,
  title={\textsc{À retenir}\ifx\relax#1\relax\relax\else\textmd{ : \itshape #1}\fi}}
\newtcolorbox{methode}[1][]{boxbase,colback=methcol!7,colframe=methcol,sharp corners,
  title={\textsc{Méthode}\ifx\relax#1\relax\relax\else\textmd{ --- \itshape #1}\fi}}
\newtcolorbox{exemplebox}[1][]{boxbase,colback=excol!5,colframe=excol,
  title={\textsc{Exemple}\ifx\relax#1\relax\relax\else\textmd{ : \itshape #1}\fi}}
\newtcolorbox{piege}[1][]{boxbase,colback=attncol!6,colframe=attncol,
  title={\textsc{Piège}\ifx\relax#1\relax\relax\else\textmd{ : \itshape #1}\fi}}
\newtcolorbox{remarque}[1][]{boxbase,colback=black!4,colframe=black!45,
  title={\textsc{Remarque}\ifx\relax#1\relax\relax\else\textmd{ : \itshape #1}\fi}}
% --- boîtes EXERCICE / CORRIGÉ (noms EXACTS, auto-comptées) ---
\newtcolorbox[auto counter]{exo}[1][]{boxbase,colback=primary!4,colframe=primary,
  title={\textsc{Exercice}~\thetcbcounter\ifx\relax#1\relax\relax\else\textmd{ --- \itshape #1}\fi}}
\newtcolorbox[auto counter]{corrige}[1][]{boxbase,colback=excol!4,colframe=excol,
  title={\textsc{Corrigé}~\thetcbcounter\ifx\relax#1\relax\relax\else\textmd{ --- \itshape #1}\fi}}
\newcommand{\R}{\mathbb{R}}\newcommand{\N}{\mathbb{N}}\newcommand{\Z}{\mathbb{Z}}\newcommand{\ds}{\displaystyle}
% (un \usepackage{fancyhdr}+en-tête est possible ; il est ignoré côté web)
% =========================================================================
\begin{document}
\newcommand{\AX}[2]{\ensuremath{\mathrm{AX_{#1}E_{#2}}}}
\begin{exo}[la famille bipyramide trigonale]
Ces trois édifices ont $5$ sites de répulsion. Pour chacun, donne la notation \AX{n}{m} et la
\textbf{géométrie moléculaire}.
\begin{enumerate}[label=\alph*)]
  \item \ce{SF4} : soufre lié à $4$ \ce{F}, $1$ doublet libre.
  \item \ce{ClF3} : chlore lié à $3$ \ce{F}, $2$ doublets libres.
  \item \ce{XeF2} : xénon lié à $2$ \ce{F}, $3$ doublets libres.
\end{enumerate}
\end{exo}

\begin{exo}[la famille octaédrique]
Ces trois édifices ont $6$ sites de répulsion. Donne \AX{n}{m} et la géométrie moléculaire.
\begin{enumerate}[label=\alph*)]
  \item \ce{SF6} : soufre lié à $6$ \ce{F}, $0$ doublet libre.
  \item \ce{BrF5} : brome lié à $5$ \ce{F}, $1$ doublet libre.
  \item \ce{XeF4} : xénon lié à $4$ \ce{F}, $2$ doublets libres.
\end{enumerate}
\end{exo}

\begin{exo}[pourquoi le doublet va-t-il à l'équateur ?]
Dans \ce{SF4} (\AX{4}{1}), les $5$ sites forment une bipyramide trigonale : $2$ positions
\emph{axiales} (à $180^\circ$) et $3$ positions \emph{équatoriales} (à $120^\circ$).
\begin{enumerate}[label=\alph*)]
  \item En position \emph{axiale}, combien de voisins à $90^\circ$ un doublet aurait-il ? Et en
        position \emph{équatoriale} ?
  \item Où se place donc le doublet libre, et quelle forme (« à bascule ») en résulte-t-il ?
\end{enumerate}
\end{exo}

\begin{exo}[la fermeture progressive de l'angle]
On compare \ce{CH4}, \ce{NH3} et \ce{H2O}, qui ont tous $4$ sites de répulsion.
\begin{enumerate}[label=\alph*)]
  \item Donne \AX{n}{m} et l'angle mesuré pour chacun ($109{,}5^\circ$, $107^\circ$, $104{,}5^\circ$
        à attribuer).
  \item Explique, avec la hiérarchie des répulsions, pourquoi l'angle diminue de \ce{CH4} à \ce{H2O}.
\end{enumerate}
\end{exo}

\begin{exo}[deux molécules linéaires… pour deux raisons opposées]
Détermine \AX{n}{m} et la géométrie moléculaire de chacun de ces édifices.
\begin{enumerate}[label=\alph*)]
  \item \ce{BeCl2} : béryllium lié à $2$ \ce{Cl}, $0$ doublet libre.
  \item \ce{SnCl2} : étain lié à $2$ \ce{Cl}, $1$ doublet libre.
  \item \ce{ICl2-} : iode lié à $2$ \ce{Cl}, $3$ doublets libres.
  \item \ce{PCl6-} : phosphore lié à $6$ \ce{Cl}, $0$ doublet libre.
  \item \ce{BeCl2} et \ce{ICl2-} sont tous deux \textbf{linéaires} : explique pourquoi, alors que
        leurs notations \AX{n}{m} sont différentes.
\end{enumerate}
\end{exo}

\end{document}
